\documentclass[]{article}
\usepackage[utf8]{inputenc}
\usepackage{amsmath, amssymb, amsthm}
\usepackage{geometry}
\usepackage{enumitem}
\usepackage{xcolor}
\usepackage[most]{tcolorbox}
\usepackage{fancyhdr}

\geometry{margin=1in}
\pagestyle{fancy}
\fancyhf{}
\rhead{AMC12/COMC Problem Analysis}
\lhead{\today}

% Color definitions
\definecolor{problemconfig}{RGB}{230, 242, 255}    % Light blue
\definecolor{strategic}{RGB}{255, 230, 242}       % Light pink
\definecolor{tactical}{RGB}{242, 255, 230}        % Light green
\definecolor{techniques}{RGB}{255, 242, 230}      % Light yellow
\definecolor{verification}{RGB}{255, 230, 230}    % Light red
\definecolor{solution}{RGB}{230, 255, 255}        % Light cyan
\definecolor{competition}{RGB}{242, 242, 255}     % Light purple
\definecolor{gold}{RGB}{218, 165, 32}

% Box styles
\newtcolorbox{problemconfigbox}{
    colback=problemconfig,
    colframe=blue,
    title=Problem Configuration Analysis,
    fonttitle=\bfseries,
    arc=3pt,
    boxrule=1pt,
    breakable
}

\newtcolorbox{strategicbox}{
    colback=strategic,
    colframe=purple,
    title=Strategic Framework Application,
    fonttitle=\bfseries,
    arc=3pt,
    boxrule=1pt,
    breakable
}

\newtcolorbox{tacticalbox}{
    colback=tactical,
    colframe=green,
    title=Tactical Methodology Selection,
    fonttitle=\bfseries,
    arc=3pt,
    boxrule=1pt,
    breakable
}

\newtcolorbox{techniquesbox}{
    colback=techniques,
    colframe=orange,
    title=Conversion Techniques Application,
    fonttitle=\bfseries,
    arc=3pt,
    boxrule=1pt,
    breakable
}

\newtcolorbox{verificationbox}{
    colback=verification,
    colframe=red,
    title=Verification Strategy Design,
    fonttitle=\bfseries,
    arc=3pt,
    boxrule=1pt,
    breakable
}

\newtcolorbox{solutionbox}{
    colback=solution,
    colframe=cyan,
    title=Solution Roadmap,
    fonttitle=\bfseries,
    arc=3pt,
    boxrule=1pt,
    breakable
}

\newtcolorbox{competitionbox}{
    colback=competition,
    colframe=purple,
    title=Competition Strategy Integration,
    fonttitle=\bfseries,
    arc=3pt,
    boxrule=1pt,
    breakable
}

\newtcolorbox{keyinsightbox}{
    colback=yellow!20,
    colframe=gold,
    title=Key Insight,
    fonttitle=\bfseries,
    arc=3pt,
    boxrule=2pt,
    leftrule=5pt,
    breakable
}

\newtcolorbox{warningbox}{
    colback=red!10,
    colframe=red!50,
    title=Warning: Common Pitfall,
    fonttitle=\bfseries,
    arc=3pt,
    boxrule=2pt,
    leftrule=5pt,
    breakable
}

\newtcolorbox{tipbox}{
    colback=green!10,
    colframe=green!50,
    title=Pro Tip,
    fonttitle=\bfseries,
    arc=3pt,
    boxrule=2pt,
    leftrule=5pt,
    breakable
}

\begin{document}

\title{AMC12B 2016 Problem 10: Strategic Analysis}
\author{Competition Math Framework}
\date{\today}
\maketitle

\section*{Problem Statement}
A quadrilateral has vertices $P(a,b)$, $Q(b,a)$, $R(-a, -b)$, and $S(-b, -a)$, where $a$ and $b$ are integers with $a>b>0$. The area of $PQRS$ is $16$. What is $a+b$?

$$\textbf{(A)}\ 4 \qquad\textbf{(B)}\ 5 \qquad\textbf{(C)}\ 6 \qquad\textbf{(D)}\ 12 \qquad\textbf{(E)}\ 13$$

\begin{problemconfigbox}
\subsection*{A. Problem Classification}
\begin{itemize}[leftmargin=*]
    \item \textbf{Problem Type}: To Find - the value of $a+b$ given geometric and algebraic constraints
    \item \textbf{Difficulty Band}: Mid-band (Problem 10/25) - requires recognizing shape properties and algebraic manipulation
    \item \textbf{Competition Context}: AMC12 (75 min, 25 problems) - approximately 3 minutes allocated per problem
    \item \textbf{Mathematical Domain}: Coordinate Geometry + Algebra, with potential Number Theory constraint satisfaction
\end{itemize}

\subsection*{B. Problem Structure Identification}
\begin{itemize}[leftmargin=*]
    \item \textbf{Given Information}: 
    \begin{itemize}
        \item Four vertices: $P(a,b)$, $Q(b,a)$, $R(-a,-b)$, $S(-b,-a)$
        \item Constraint: $a > b > 0$ and $a,b \in \mathbb{Z}$
        \item Area of quadrilateral $PQRS = 16$
    \end{itemize}
    
    \item \textbf{Constraints}: 
    \begin{itemize}
        \item Integer values only: $a, b \in \mathbb{Z}^+$
        \item Ordering constraint: $a > b > 0$
        \item Area constraint: $\text{Area}(PQRS) = 16$
    \end{itemize}
    
    \item \textbf{Objective}: Find the value of $a+b$
    
    \item \textbf{Hidden Assumptions}: 
    \begin{itemize}
        \item The quadrilateral is convex (based on vertex arrangement)
        \item The vertices form a specific symmetric pattern
        \item The symmetry suggests special properties (potentially a rectangle or rhombus)
    \end{itemize}
    
    \item \textbf{Answer Choice Leverage}: 
    \begin{itemize}
        \item Small values (4, 5, 6) vs. larger values (12, 13)
        \item If area = 16 and we need integer solutions, possible factor pairs are limited
        \item Can work backwards: test answer choices to see which produces area of 16
    \end{itemize}
\end{itemize}

\subsection*{C. Strategic Orientation}
\begin{itemize}[leftmargin=*]
    \item \textbf{Hypothesis}: The symmetric coordinate structure suggests the quadrilateral has special properties
    
    \item \textbf{Conclusion}: Need to find $a+b$ where $a,b$ are positive integers satisfying the area constraint
    
    \item \textbf{Notation}: 
    \begin{itemize}
        \item $P = (a,b)$, $Q = (b,a)$, $R = (-a,-b)$, $S = (-b,-a)$
        \item Area formula to be determined based on shape recognition
    \end{itemize}
    
    \item \textbf{Intuitive Approach}: 
    \begin{itemize}
        \item Recognize the symmetry: vertices are reflections across both axes and the origin
        \item This suggests the quadrilateral might be a rectangle or special parallelogram
        \item Use Shoelace Theorem or recognize shape properties to find area formula
    \end{itemize}
    
    \item \textbf{Cross-Domain Potential}: 
    \begin{itemize}
        \item Pure coordinate geometry approach (Shoelace Theorem)
        \item Vector/linear algebra approach (diagonals, perpendicularity)
        \item Algebraic constraint solving (factor $a^2 - b^2 = 8$)
    \end{itemize}
\end{itemize}
\end{problemconfigbox}

\begin{strategicbox}
\subsection*{A. Psychological Strategies}
\begin{itemize}[leftmargin=*]
    \item \textbf{Mental Toughness}: If initial approach fails, have backup: (1) Shoelace formula, (2) prove it's a rectangle, (3) answer choice elimination
    
    \item \textbf{Creativity Cultivation}: The symmetry of coordinates is a gift - exploit it! Consider: what shapes have this exact symmetry pattern?
    
    \item \textbf{Iterative Refinement}: 
    \begin{itemize}
        \item First pass: recognize symmetry
        \item Second pass: identify shape (rectangle)
        \item Third pass: calculate area using shape properties
        \item Fourth pass: solve algebraic constraint
    \end{itemize}
\end{itemize}

\subsection*{B. Investigation Strategies}
\begin{itemize}[leftmargin=*]
    \item \textbf{Startup Orientation}: This is a coordinate geometry problem with algebraic constraints - looks accessible for mid-band
    
    \item \textbf{Get Your Hands Dirty}: 
    \begin{itemize}
        \item Plot the four points on coordinate plane
        \item Observe: $P$ and $R$ are opposite through origin, same with $Q$ and $S$
        \item Calculate slopes of adjacent sides to check for perpendicularity
    \end{itemize}
    
    \item \textbf{Penultimate Step}: The final step will be computing $a+b$ once we solve for $a$ and $b$ individually
    
    \item \textbf{Wishful Thinking}: "I wish this were a rectangle!" - and it probably is based on symmetry
    
    \item \textbf{Make It Easier}: 
    \begin{itemize}
        \item Try $a=3, b=1$: vertices become $P(3,1), Q(1,3), R(-3,-1), S(-1,-3)$
        \item Check if this gives area 16
    \end{itemize}
    
    \item \textbf{Conjecture via Small Cases}: 
    \begin{itemize}
        \item $a=2, b=1$: Check area
        \item $a=3, b=1$: Check area
        \item $a=3, b=2$: Check area
    \end{itemize}
    
    \item \textbf{Visualization}: 
    \begin{itemize}
        \item Draw coordinate plane with all four points
        \item Connect vertices in order: $P \to Q \to R \to S \to P$
        \item Observe symmetry about both coordinate axes
    \end{itemize}
\end{itemize}

\subsection*{C. Late-Band Strategic Playbook}
\textit{(Not strictly late-band, but techniques still applicable)}
\begin{itemize}[leftmargin=*]
    \item \textbf{Answer-Choice Leverage}: 
    \begin{itemize}
        \item For answer (A): $a+b=4$ with $a>b>0$, possible: $(3,1)$ only (since $(2,2)$ violates $a>b$)
        \item So if (A), then $a=3, b=1$ - verify area
    \end{itemize}
    
    \item \textbf{Make It Easier}: Consider special case where $b=1$, solve for $a$
    
    \item \textbf{Penultimate Step}: Before computing $a+b$, need to solve $a^2 - b^2 = 8$ with integer constraints
\end{itemize}

\subsection*{D. Argument Construction Strategy}
\begin{itemize}[leftmargin=*]
    \item \textbf{Direct Proof Path}:
    \begin{enumerate}
        \item Prove $PQRS$ is a rectangle by showing perpendicular adjacent sides
        \item Express area as product of side lengths
        \item Set equal to 16 and solve algebraically
        \item Apply integer constraint to find unique solution
    \end{enumerate}
    
    \item \textbf{Alternative: Shoelace Approach}:
    \begin{enumerate}
        \item Apply Shoelace formula directly
        \item Simplify to get $2a^2 - 2b^2 = 16$
        \item Factor as $a^2 - b^2 = 8$
        \item Solve Diophantine equation
    \end{enumerate}
\end{itemize}
\end{strategicbox}

\begin{tacticalbox}
\subsection*{A. Core Mathematical Tactics}
\begin{itemize}[leftmargin=*]
    \item \textbf{Symmetry Exploitation}: 
    \begin{itemize}
        \item The vertices exhibit both origin symmetry and axis symmetry
        \item $P$ and $R$ are opposite through origin: $R = -P$
        \item $Q$ and $S$ are opposite through origin: $S = -Q$
        \item $P(a,b)$ and $Q(b,a)$ are reflections across line $y=x$
    \end{itemize}
    
    \item \textbf{Extremal Principle}: Not directly applicable here
    
    \item \textbf{Invariant Identification}: The symmetry pattern is preserved under reflection
    
    \item \textbf{Small Cases Analysis}:
    \begin{itemize}
        \item Test $(a,b) = (3,1)$: vertices at $(3,1), (1,3), (-3,-1), (-1,-3)$
        \item Test $(a,b) = (4,2)$: check if area = 16
        \item Test $(a,b) = (5,3)$: check if area = 16
    \end{itemize}
    
    \item \textbf{Algebraic Manipulation}:
    \begin{itemize}
        \item Factor difference of squares: $a^2 - b^2 = (a-b)(a+b)$
        \item Given $a^2 - b^2 = 8$, find factor pairs of 8
        \item Possible: $(a-b, a+b) = (1,8), (2,4), (4,2), (8,1)$
    \end{itemize}
    
    \item \textbf{Coordinate Geometry Tools}:
    \begin{itemize}
        \item Slope formula: $m_{PQ} = \frac{a-b}{b-a} = -1$
        \item Distance formula: $|PQ| = \sqrt{(a-b)^2 + (b-a)^2} = \sqrt{2}|a-b|$
        \item Shoelace Theorem for area calculation
    \end{itemize}
\end{itemize}

\subsection*{B. Crossover Tactics}
\begin{itemize}[leftmargin=*]
    \item \textbf{Graph Theory}: Not applicable
    
    \item \textbf{Complex Numbers}: 
    \begin{itemize}
        \item Could represent points as $P = a + bi$, $Q = b + ai$, etc.
        \item Rotation by $90^\circ$ in complex plane: multiply by $i$
        \item But this approach adds unnecessary complexity
    \end{itemize}
    
    \item \textbf{Linear Algebra}:
    \begin{itemize}
        \item View as vectors: $\vec{PQ} = (b-a, a-b)$
        \item Check perpendicularity via dot product: $\vec{PQ} \cdot \vec{PS} = 0$
    \end{itemize}
    
    \item \textbf{Number Theory}:
    \begin{itemize}
        \item Solve Diophantine equation: $(a-b)(a+b) = 8$
        \item Find positive integer solutions with $a > b > 0$
        \item Factor pairs of 8: $(1,8), (2,4), (4,2), (8,1)$
    \end{itemize}
\end{itemize}
\end{tacticalbox}

\begin{techniquesbox}
\subsection*{A. Recognition Index Analysis}

\textbf{High-Probability Techniques Identified:}
\begin{enumerate}[leftmargin=*]
    \item \textbf{Coordinate Bash (90\% relevance)}: Direct calculation using coordinates
    \begin{itemize}
        \item Trigger: Four explicit coordinate points given
        \item Application: Use Shoelace formula or distance/perpendicularity checks
    \end{itemize}
    
    \item \textbf{Symmetry Exploitation (95\% relevance)}: Central to this problem
    \begin{itemize}
        \item Trigger: Vertices show clear pattern $(a,b) \leftrightarrow (b,a)$ and sign flips
        \item Application: Recognize rectangle/special quadrilateral from symmetry
    \end{itemize}
    
    \item \textbf{Algebraic Constraint Solving (85\% relevance)}: 
    \begin{itemize}
        \item Trigger: Area constraint + integer constraints
        \item Application: Set up equation, factor, solve for integers
    \end{itemize}
    
    \item \textbf{Difference of Squares Factorization (80\% relevance)}:
    \begin{itemize}
        \item Trigger: Expression $a^2 - b^2$ appears naturally
        \item Application: Factor as $(a-b)(a+b) = 8$ to find solutions
    \end{itemize}
    
    \item \textbf{Answer Choice Elimination (70\% relevance)}:
    \begin{itemize}
        \item Trigger: Multiple choice with small number of options
        \item Application: Test $a+b = 4$ directly by finding valid $(a,b)$ pairs
    \end{itemize}
\end{enumerate}

\subsection*{B. Technique Selection Rationale}

\textbf{Primary Technique}: Shape Recognition + Algebraic Constraint
\begin{itemize}[leftmargin=*]
    \item \textbf{Why}: The symmetry immediately suggests a rectangle, which has simple area formula
    \item \textbf{Efficiency}: Avoids tedious Shoelace calculations
    \item \textbf{Reliability}: Geometric insight leads to clean algebraic equation
\end{itemize}

\textbf{Secondary Technique}: Shoelace Theorem
\begin{itemize}[leftmargin=*]
    \item \textbf{Why}: Guaranteed to work for any quadrilateral
    \item \textbf{Trade-off}: More computation but no insight required
    \item \textbf{Use case}: Backup if shape recognition fails
\end{itemize}

\textbf{Tertiary Technique}: Answer Testing
\begin{itemize}[leftmargin=*]
    \item \textbf{Why}: Fast elimination in time-pressured scenario
    \item \textbf{Trade-off}: Less rigorous but competition-efficient
    \item \textbf{Use case}: Under time pressure or for verification
\end{itemize}

\subsection*{C. Integration Strategy}

\textbf{Combined Approach Flow:}
\begin{enumerate}
    \item \textbf{Step 1 - Recognition (30 sec)}: Identify symmetry pattern, hypothesize rectangle
    \item \textbf{Step 2 - Verification (45 sec)}: Prove perpendicularity of adjacent sides
    \item \textbf{Step 3 - Area Formula (30 sec)}: Calculate side lengths, express area
    \item \textbf{Step 4 - Algebraic Solve (60 sec)}: Set area = 16, simplify to $a^2-b^2=8$
    \item \textbf{Step 5 - Factor and Solve (45 sec)}: Factor as $(a-b)(a+b)=8$, find integer solutions
\end{enumerate}

\textbf{Conflict Resolution:}
\begin{itemize}[leftmargin=*]
    \item If rectangle proof seems difficult, pivot to Shoelace immediately
    \item If algebra becomes messy, test answer choices directly
    \item If multiple solutions appear, apply constraint $a > b > 0$ to eliminate
\end{itemize}
\end{techniquesbox}

\begin{verificationbox}
\subsection*{A. Verification Level Selection}

\textbf{Decision Tree Analysis:}
\begin{itemize}[leftmargin=*]
    \item Problem difficulty: Mid-band (Problem 10) → suggests Level 3-4 verification
    \item Multiple approaches available → confidence can be high with dual verification
    \item Integer constraint provides strong check → use divisibility/factorization tests
\end{itemize}

\textbf{Selected Level}: \textbf{Level 4 - Comprehensive Verification}
\begin{itemize}[leftmargin=*]
    \item Reasoning: Mid-band problem with clean answer suggests thorough check is efficient
    \item Time allocated: 60 seconds (within 3-minute total budget)
\end{itemize}

\subsection*{B. Verification Methods}

\textbf{Primary Verification - Direct Substitution:}
\begin{enumerate}
    \item For $a=3, b=1$:
    \begin{itemize}
        \item Vertices: $P(3,1), Q(1,3), R(-3,-1), S(-1,-3)$
        \item Check area = 16 using rectangle formula or Shoelace
        \item Verify $a > b > 0$: $3 > 1 > 0$ $\checkmark$
    \end{itemize}
\end{enumerate}

\textbf{Secondary Verification - Alternative Method:}
\begin{enumerate}
    \item Use Shoelace formula independently:
    \begin{align*}
    \text{Area} &= \frac{1}{2}|x_1(y_2-y_4) + x_2(y_3-y_1) + x_3(y_4-y_2) + x_4(y_1-y_3)|\\
    &= \frac{1}{2}|a(a-(-a)) + b(-b-b) + (-a)(-a-a) + (-b)(b-(-b))|\\
    &= \frac{1}{2}|2a^2 - 2b^2 + 2a^2 - 2b^2|\\
    &= 2a^2 - 2b^2 = 2(a^2-b^2)
    \end{align*}
    \item Setting equal to 16: $2(a^2-b^2) = 16 \Rightarrow a^2-b^2=8$
    \item For $a=3, b=1$: $9-1=8$ $\checkmark$
\end{enumerate}

\textbf{Tertiary Verification - Uniqueness Check:}
\begin{itemize}[leftmargin=*]
    \item Factor pairs of 8: $(1,8), (2,4), (4,2), (8,1)$
    \item For $(a-b)(a+b)=8$:
    \begin{itemize}
        \item $(a-b, a+b) = (1,8)$: $a=4.5$ (not integer) $\times$
        \item $(a-b, a+b) = (2,4)$: $a=3, b=1$ $\checkmark$
        \item $(a-b, a+b) = (4,2)$: $a=3, b=-1$ (violates $b>0$) $\times$
        \item $(a-b, a+b) = (8,1)$: $a=4.5$ (not integer) $\times$
    \end{itemize}
    \item Only solution: $(a,b) = (3,1)$, so $a+b=4$ $\checkmark$
\end{itemize}

\subsection*{C. Confidence Calibration}

\textbf{Confidence Assessment}: \textbf{98\%}
\begin{itemize}[leftmargin=*]
    \item Three independent verifications all confirm answer
    \item Algebraic factorization shows uniqueness of solution
    \item Geometric interpretation (rectangle) is consistent
    \item Answer matches option (A) from multiple choice
\end{itemize}

\textbf{Remaining 2\% doubt sources:}
\begin{itemize}[leftmargin=*]
    \item Potential computational error in Shoelace formula
    \item Possible misreading of constraint $a>b>0$
    \item Could have missed a factor pair (though checked systematically)
\end{itemize}

\textbf{Flag for Review}: \textbf{No} - confidence is very high and time is better spent on harder problems
\end{verificationbox}

\begin{solutionbox}
\subsection*{A. Step-by-Step Solution Roadmap}

\textbf{Solution Path 1: Shape Recognition + Algebra (Primary, 3 minutes)}

\textbf{Phase 1 - Shape Identification (60 seconds)}
\begin{enumerate}
    \item Plot the four vertices and observe symmetry pattern
    \item Calculate slope of $PQ$: 
    $$m_{PQ} = \frac{a-b}{b-a} = -1$$
    \item Calculate slope of $PS$: 
    $$m_{PS} = \frac{-a-b}{-b-a} = \frac{-(a+b)}{-(b+a)} = 1$$
    \item Since $m_{PQ} \cdot m_{PS} = (-1)(1) = -1$, we have $PQ \perp PS$
    \item \textit{Checkpoint}: Confirmed $PQRS$ is a rectangle (or at least has one right angle)
\end{enumerate}

\textbf{Phase 2 - Area Calculation (60 seconds)}
\begin{enumerate}
    \item Calculate length $|PQ|$:
    $$|PQ| = \sqrt{(b-a)^2 + (a-b)^2} = \sqrt{2(a-b)^2} = \sqrt{2}|a-b|$$
    \item Calculate length $|PS|$:
    $$|PS| = \sqrt{(-b-a)^2 + (-a-b)^2} = \sqrt{2(a+b)^2} = \sqrt{2}|a+b|$$
    \item Area of rectangle:
    $$\text{Area} = |PQ| \times |PS| = \sqrt{2}(a-b) \times \sqrt{2}(a+b) = 2(a-b)(a+b)$$
    \item \textit{Checkpoint}: Area formula derived as $2(a^2-b^2)$
\end{enumerate}

\textbf{Phase 3 - Algebraic Solution (60 seconds)}
\begin{enumerate}
    \item Set area equal to 16:
    $$2(a^2-b^2) = 16$$
    $$a^2 - b^2 = 8$$
    \item Factor:
    $$(a-b)(a+b) = 8$$
    \item Find factor pairs of 8: $(1,8), (2,4), (4,2), (8,1)$
    \item Test each with constraints $a,b \in \mathbb{Z}^+$ and $a>b$:
    \begin{itemize}
        \item $(a-b, a+b) = (2,4)$: solving gives $a=3, b=1$ $\checkmark$
    \end{itemize}
    \item \textit{Checkpoint}: Found unique solution $(a,b) = (3,1)$
\end{enumerate}

\textbf{Phase 4 - Final Answer (10 seconds)}
\begin{enumerate}
    \item Compute: $a + b = 3 + 1 = 4$
    \item Select answer choice: $\boxed{\textbf{(A)}\ 4}$
\end{enumerate}

\subsection*{B. Alternative Approaches}

\textbf{Solution Path 2: Shoelace Theorem (Backup, 4 minutes)}
\begin{enumerate}
    \item Apply Shoelace formula with vertices in order $P,Q,R,S$
    \item Substitute coordinates and simplify to get $2(a^2-b^2) = 16$
    \item Proceed with same algebra as Path 1
\end{enumerate}

\textbf{Solution Path 3: Answer Testing (Competition shortcut, 2 minutes)}
\begin{enumerate}
    \item Test answer (A): $a+b=4$ with $a>b>0$ gives $(a,b)=(3,1)$
    \item For $(3,1)$: compute area using any method
    \item If area = 16, done. If not, try next answer choice.
\end{enumerate}

\subsection*{C. Pitfall Guide}

\textbf{Critical Pitfalls:}
\begin{enumerate}
    \item \textbf{Vertex Order Error}: 
    \begin{itemize}
        \item Mistake: Connecting vertices in wrong order
        \item Prevention: Plot points first, connect in sequence $P \to Q \to R \to S$
    \end{itemize}
    
    \item \textbf{Shoelace Sign Error}:
    \begin{itemize}
        \item Mistake: Forgetting absolute value in Shoelace formula
        \item Prevention: Always take absolute value of final Shoelace expression
    \end{itemize}
    
    \item \textbf{Incomplete Factorization}:
    \begin{itemize}
        \item Prevention: Systematically list all factor pairs including reversed order
    \end{itemize}
    
    \item \textbf{Constraint Violation}:
    \begin{itemize}
        \item Mistake: Accepting $(a,b) = (3,-1)$ from $(4,2)$ factor pair
        \item Prevention: Check all three constraints: $a,b \in \mathbb{Z}^+$ and $a>b$
    \end{itemize}
\end{enumerate}

\subsection*{D. Time Management}

\textbf{Time Allocation:}
\begin{itemize}[leftmargin=*]
    \item Initial reading: 15 seconds
    \item Shape recognition: 45 seconds
    \item Area calculation: 60 seconds  
    \item Algebraic solving: 60 seconds
    \item Verification: 30 seconds
    \item \textbf{Total: 3.5 minutes}
\end{itemize}

\textbf{Pivot Indicators:}
\begin{itemize}[leftmargin=*]
    \item If perpendicularity check takes >2 minutes → switch to Shoelace
    \item If algebra yields no integer solutions after 1 minute → recheck area calculation
    \item If total time exceeds 5 minutes → mark answer, move on
\end{itemize}
\end{solutionbox}

\begin{competitionbox}
\subsection*{A. Time Management Strategy}

\textbf{Problem 10 Time Budget:}
\begin{itemize}[leftmargin=*]
    \item Standard allocation: 3 minutes
    \item Recommended for P10: 3-4 minutes (mid-band problems deserve extra time)
    \item Maximum acceptable: 5 minutes before abandonment consideration
\end{itemize}

\subsection*{B. Problem Prioritization}

\textbf{Priority Assessment for P10:}
\begin{itemize}[leftmargin=*]
    \item \textbf{Accessibility}: High - coordinate geometry with visible symmetry
    \item \textbf{Familiarity}: Medium-High - standard rectangle area + constraint solving
    \item \textbf{Time Efficiency}: Good - clear path to solution in 3-4 minutes
    \item \textbf{Confidence Potential}: High - multiple verification methods available
\end{itemize}

\subsection*{C. Risk Management}

\textbf{Confidence Assessment Protocol:}
\begin{itemize}[leftmargin=*]
    \item After solving, current confidence: 95\%+
    \item Verification level: Comprehensive (three methods)
    \item Decision: Mark answer, do not flag for review
\end{itemize}

\subsection*{D. Answer Format Management}

\textbf{AMC12 Multiple Choice Protocol:}
\begin{itemize}[leftmargin=*]
    \item Final answer: $a+b = 4$
    \item Corresponding choice: \textbf{(A)}
    \item Bubble carefully on answer sheet
    \item Double-check problem number 10 matches answer line 10
\end{itemize}
\end{competitionbox}

\begin{keyinsightbox}
\textbf{Key Insight}: The symmetry of the coordinates immediately reveals that $PQRS$ is a rectangle. The vertices $P(a,b)$ and $Q(b,a)$ are reflections across $y=x$, while $R(-a,-b)$ and $S(-b,-a)$ are their opposites through the origin. This symmetry guarantees perpendicular sides.

Once recognized as a rectangle, the area simplifies to $2(a-b)(a+b) = 2(a^2-b^2)$, leading directly to the factorization $(a-b)(a+b) = 8$. The integer constraint with $a>b>0$ uniquely determines $(a,b) = (3,1)$.
\end{keyinsightbox}

\begin{warningbox}
\textbf{Warning}: The most common error is incorrectly computing the area using Shoelace without properly ordering the vertices. Always connect $P \to Q \to R \to S \to P$ in sequence. Another frequent mistake is forgetting to check all factor pairs of 8 or missing the constraint $b>0$ when testing $(4,2)$ which gives $b=-1$.
\end{warningbox}

\begin{tipbox}
\textbf{Pro Tip}: In a time-pressured competition setting, immediately test answer choice (A) by assuming $a+b=4$. With $a>b>0$ and both integers, try $(a,b)=(3,1)$. Quickly compute the area. If it equals 16, you're done in under 2 minutes!

For verification, the factorization approach is fastest: $a^2-b^2=8$ factors as $(a-b)(a+b)=8$. Testing $(a-b, a+b) = (2,4)$ immediately gives $(a,b)=(3,1)$ and confirms $a+b=4$.
\end{tipbox}

\section*{Final Answer}
$$\boxed{\textbf{(A)}\ 4}$$

The unique solution is $(a,b) = (3,1)$, giving $a+b=4$.

\section*{Confidence Assessment}
\begin{itemize}[leftmargin=*]
    \item \textbf{Confidence Level}: 98\% - Three independent verification methods confirm the answer
    \item \textbf{Verification Applied}: Direct substitution, Shoelace formula, and factor pair uniqueness analysis
    \item \textbf{Flag for Review}: No - extremely high confidence
    \item \textbf{Time Spent}: 3.5 minutes actual vs. 3 minutes estimated
\end{itemize}

\end{document}